\section{模型评价与推广}\label{sec:eval}

\paragraph{优点.}
(i) 四问共用同一套散度形式非线性方程与节点型有限体积离散，物理主线一致、可追溯；
(ii) 达标判据在未舍入全域重构场上连续求根，区分阈值穿越时间 $\tstar$ 与严格合格采样时刻，
避免四位小数显示的伪精度；(iii) 问题四以参考坐标处理动边界，方程只含 $R(t)$，守恒离散精确；
(iv) 数值配置由各问实际场点误差选定并登记，界面系数、时间格式、网格均通过独立检验。

\paragraph{局限：端面效应（V-15 增量校核）.}
一维径向近似忽略端面。以线性常系数、恒环境阶跃辅助问题（乘积解 $\theta_{\mathrm{finite}}=\theta_{\mathrm{radial}}
\theta_{\mathrm{slab}}$，平板半厚 $L/2=\SI{0.125}{m}$）估计端面\textbf{增量} $\Delta T=(T_0-T_\infty)
\theta_{\mathrm{radial}}(\theta_{\mathrm{slab}}-1)$（质同理），在问题一至四\emph{全部相关表点时刻/位置}
（含问题二表 3/4 的 $0.5$--$\SI{3}{h}$）评估：最大 $|\Delta T|=\SI{9.1013e-4}{\celsius}$（位于 $\SI{1.5}{h}$、$r=0$）、
$|\Delta C|=\num{1.1281e-5}$（位于 $\SI{42}{h}$、$r=0$），均远低于 $\SI{0.1}{\celsius}$、$0.005$。
$|\Delta C|$ 随平板本征项数 $\text{nterms}=60/120/200$ 收敛于 $\num{1.13e-5}$（$60$ 项时的 $\num{4.76e-4}$
系小 $Fo$ 级数截断底噪，非物理值），故报告值为物理增量。\textbf{该结论仅限固定物性、固定环境的线性辅助
模型及已检查位置/时刻，不认证/否定变物性时变环境的非线性主线。} 机理：轴向扩散时间尺度
$(L/2)^2/D\approx\SI{886}{h}$ 远大于径向 $R_0^2/D\approx\SI{22.7}{h}$，端面显现时径向场已趋均衡，
故乘积增量极小。

\paragraph{其他局限.}
本文采用题给有效传热传质模型，未显式解析相变潜热；更完整机理模型需要一致的界面平衡及材料参数，
超出本题给定信息。烘房「水分浓度」作为有效边界驱动 $C_{\mathrm{env}}$，非已识别的真实平衡含水率；
$h,h_m$ 跨问沿用为额外有效系数假设。$\SI{4}{h}$ 后环境与 $\SI{72}{h}$ 后半径均为外推，已显式标注。

\paragraph{推广.}
模型框架可推广到其他细长物料的对流干燥：更换经验物性组、边界系数与几何驱动即可；
积分界面系数对陡峭扩散前沿的稳健性、参考坐标对给定收缩轨迹的处理，适用于一类含强非线性扩散与
动边界的干燥/浸出问题。若获得材料吸附--解吸等温线与相变参数，可在现框架上增补平衡边界与潜热汇。
